% T I T L E   P A G E
% -------------------
\pagestyle{empty}
\pagenumbering{roman}

\begin{titlepage}
    \begin{center}
    \normalsize
    \MakeUppercase{University of California, Los Angeles} \\
    Department of Mechanical and Aerospace Engineering \\

    \vspace{3.0cm}

    \Large
    {\bf Optimal Control of \\ MRI-Guided Microwave Thermal Ablation} \\

    \vspace{4.0cm}

    \normalsize
    MAE 270C — Optimal Control Final Project Report \\

    \vspace{2.5cm}

    by \\
    \vspace{1.0cm}
    Jing-Yuan (Kevin) Huang

    \vspace{2.0cm}

    Advisor: Professor Shahriar Talebi \\
    \vspace{0.5cm}
    2026
    \end{center}
\end{titlepage}

\pagestyle{plain}
\setcounter{page}{2}

\cleardoublepage
\phantomsection

% A B S T R A C T
% ---------------
\addcontentsline{toc}{chapter}{Abstract}
\begin{center}\textbf{Abstract}\end{center}

\noindent
This report presents a PDE-constrained optimal control framework for
MRI-guided microwave thermal ablation of liver tumors.
The state equations consist of the Pennes bioheat partial differential equation
(discretized to a finite-difference ODE) and an Arrhenius thermal damage ODE.
The control objective is formulated as a Bolza problem with free final time:
complete tumor ablation (Arrhenius damage $\Omega_d \geq 1$ in $\Omega_T$)
is penalized at the terminal time, while healthy-tissue overheating,
total energy, and treatment duration are penalized in the running cost.

Three solution strategies are implemented and compared:
(i)~open-loop bang-bang simulation,
(ii)~indirect Pontryagin-based gradient projection (forward-adjoint sweeps), and
(iii)~direct single-shooting nonlinear programming via \texttt{scipy} SLSQP.
The framework supports both 2-D slices and full 3-D volumes, three probe
(SAR) geometries, and Dirichlet, Neumann, and Robin boundary conditions.
A sweep over solver, dimensionality, and probe geometry demonstrates the
trade-off between treatment efficacy and collateral thermal damage, and shows
how cost-weight calibration and solver scalability behave as the problem moves
from two to three dimensions.

\cleardoublepage
\phantomsection

% A C K N O W L E D G E M E N T S
% --------------------------------
\addcontentsline{toc}{chapter}{Acknowledgements}
\begin{center}\textbf{Acknowledgements}\end{center}

\noindent
Huge thanks to Professor Shahriar Talebi for his wonderful teaching and mentorship throughout MAE 270C.
This is one of the best control courses I have ever taken.

Also thanks to AI agent Claude for helping me construct the codes and visualization of simulation results.

\cleardoublepage
\phantomsection

% T A B L E   O F   C O N T E N T S
% ----------------------------------
\renewcommand\contentsname{Table of Contents}
\tableofcontents
\cleardoublepage
\phantomsection

% L I S T   O F   S Y M B O L S
% ------------------------------
\printglossary[type=symbols]
\cleardoublepage
\phantomsection

\pagenumbering{arabic}
